\documentclass[11pt,a4paper]{article}
\usepackage{fontspec,amsmath,amssymb,amsthm,microtype}
\setmainfont{lmroman10-regular.otf}[BoldFont=lmroman10-bold.otf,ItalicFont=lmroman10-italic.otf,BoldItalicFont=lmroman10-bolditalic.otf]
\usepackage{xcolor}
\usepackage[most]{tcolorbox}
\usepackage{mathrsfs}
\providecommand{\lt}{<}
\providecommand{\gt}{>}
\setmonofont{DejaVuSansMono.ttf}[BoldFont=DejaVuSansMono-Bold.ttf,ItalicFont=DejaVuSansMono-Oblique.ttf,BoldItalicFont=DejaVuSansMono-BoldOblique.ttf,Scale=0.86]
\newfontfamily\leanfallback{STIXTwoMath-Regular.otf}[Scale=0.86]
\definecolor{leanrule}{gray}{0.62}
\definecolor{leanback}{gray}{0.962}
\definecolor{leantitle}{gray}{0.22}
\definecolor{leanmuted}{gray}{0.38}
\newtcolorbox{leanbox}[1]{enhanced,breakable,colback=leanback,colframe=leanrule,
  boxrule=0.4pt,arc=0.6mm,left=2.4mm,right=2.4mm,top=1.8mm,bottom=1.8mm,
  before skip=0.8em,after skip=1em,fonttitle=\small\bfseries,coltitle=leantitle,
  colbacktitle=leanback,titlerule=0.3pt,title={#1}}
\newlength{\leancharwidth}
\newenvironment{leancode}{\par\smallskip\begingroup\ttfamily\footnotesize\raggedright
  \settowidth{\leancharwidth}{M}\setlength{\parskip}{0pt}\setlength{\parindent}{0pt}
  \frenchspacing}{\par\endgroup\smallskip}
\newcommand{\leanline}[2]{\hangindent=\dimexpr#1\leancharwidth+4\leancharwidth\relax
  \hangafter=1\noindent\hspace*{#1\leancharwidth}#2\par}
\newcommand{\leansource}[1]{\par{\footnotesize\color{leanmuted}#1}\par}
\newcommand{\leanlabel}[1]{\par\smallskip\noindent{\small\bfseries #1}\par}
\usepackage[margin=25mm]{geometry}
\usepackage[unicode,hidelinks]{hyperref}
\hypersetup{pdfauthor={LangevinLean Project}}
\setlength{\emergencystretch}{3em}
\setlength{\parskip}{0.4em}
\setlength{\parindent}{0pt}
\allowdisplaybreaks[2]
\raggedbottom
\title{\texorpdfstring{Euler convergence with a moving step and endpoint}{Euler convergence with a moving step and endpoint}}
\author{LangevinLean Project}
\date{}
\begin{document}\maketitle
Draft mathematical exposition; these notes may contain mistakes.


For a Lipschitz drift and continuous additive input, Euler approximations converge to the actual integral solution when the step tends to zero and the elapsed grid time tends to the desired time. The step and iteration count may vary independently along any filter subject to a fixed finite horizon bound.



Earlier results used here: \cite[Section~3.16]{LangevinLean2026ULASC}.

\section{\texorpdfstring{Uniform finite-horizon control gives the moving-time limit}{Uniform finite-horizon control gives the moving-time limit}}


\textbf{Claim.} If \(r_i\to0\), \(n_ir_i\le T\) and \(n_ir_i\to t\), then the actual additive Euler endpoint converges to \(Y_t\).



\textbf{Proof.} The deterministic finite-horizon Euler estimate bounds the endpoint error relative to \(Y_{n_ir_i}\) by a constant times the modulus of continuity of the drift along the solution over intervals of length \(r_i\). Uniform continuity on the compact horizon makes that modulus tend to zero. Continuity of \(Y\) makes \(Y_{n_ir_i}\to Y_t\). The triangle inequality combines these limits. The case of a zero Lipschitz constant is included in the same finite-horizon estimate.

\begin{leanbox}{Lean: \texttt{additiveEulerChain\_\allowbreak{}variable\_\allowbreak{}tendsto}}
\leansource{\texttt{additiveEulerChain\_\allowbreak{}variable\_\allowbreak{}tendsto} \quad WassersteinMixing/\allowbreak{}SemigroupEuler.lean, lines 22–70}
\begin{leancode}
\leanline{0}{theorem\ additiveEulerChain\_variable\_tendsto}
\leanline{4}{(g\ :\ H\ →\ H)\ (L\ :\ ℝ≥0)\ (hLip\ :\ LipschitzWith\ L\ g)}
\leanline{4}{(x\ :\ H)\ (B\ Y\ :\ ℝ\ →\ H)\ (hB0\ :\ B\ 0\ =\ 0)\ (hY\ :\ Continuous\ Y)}
\leanline{4}{(T\ :\ ℝ)\ (hT\ :\ 0\ ≤\ T)}
\leanline{4}{(hYe\ :\ ∀\ t\ ∈\ Icc\ 0\ T,\ Y\ t\ -\ x\ =\ -(∫\ s\ in\ 0..t,\ g\ (Y\ s))\ +\ B\ t)}
\leanline{4}{(step\ :\ I\ →\ ℝ)\ (count\ :\ I\ →\ ℕ)\ (hstep\ :\ ∀\ i,\ 0\ ≤\ step\ i)}
\leanline{4}{(hmesh\ :\ Tendsto\ step\ l\ ({\leanfallback 𝓝}\ 0))\ (hbound\ :\ ∀\ i,\ (count\ i\ :\ ℝ)\ *\ step\ i\ ≤\ T)}
\leanline{4}{(t\ :\ ℝ)\ (htime\ :\ Tendsto\ (fun\ i\ ↦\ (count\ i\ :\ ℝ)\ *\ step\ i)\ l\ ({\leanfallback 𝓝}\ t))\ :}
\leanline{4}{Tendsto\ (fun\ i\ ↦\ ULA.additiveEulerChain\ g\ (step\ i)\ x\ B\ (count\ i))\ l\ ({\leanfallback 𝓝}\ (Y\ t))}
\end{leancode}
\leanlabel{How the Lean proof follows the argument}
\texttt{additive\allowbreak{}Euler\allowbreak{}Chain\_\allowbreak{}variable\_\allowbreak{}tendsto} applies the proved uniform additive Euler error estimate and compact-horizon uniform continuity. It separates the Euler error at the actual elapsed grid time from the moving-endpoint continuity error, then combines their limits.
\end{leanbox}


Earlier results used here: \cite[Section~3.16.1]{LangevinLean2026ULASC}; \cite[Section~3.16.7]{LangevinLean2026ULASC}.

\begin{thebibliography}{ULA-LSI}
\bibitem[ULA-SC]{LangevinLean2026ULASC}
LangevinLean Project. \emph{Unadjusted Langevin under Strong Convexity}. LL-01, 2026. Draft report, content revision d34ead06987c. \url{https://odakyildiz.github.io/LangevinLean/docs/pdf/UnadjustedLangevin/FullProof.pdf}.
\end{thebibliography}

\end{document}
